\chapter{Production deployment és service lifecycle}

Ez a fejezet nem ismétli meg a konfigurációs referencia és az observability részleteit. A production szempontból azt foglalja össze, hogyan lesz egy ellenőrzött release-ből futó service, hogyan történik a release-váltás, és hogyan marad kontrollált a process lifecycle. A teljes backup/restore és disaster recovery a következő fejezet önálló témája.

\section{Release útvonal}

Ajánlott sorrend:

\begin{lstlisting}
./verify.sh
-> cargo build --locked --release
-> immutable release directory
-> rwlang-server --config ... --check-config
-> rwlang-cli migrate verify
-> approved rwlang-cli migrate apply
-> rwlang-cli migrate verify
-> atomic current symlink switch
-> controlled restart
-> /health/live
-> /health/ready
-> log / metric / audit review
\end{lstlisting}

A buildelt \code{rwlang-server} és \code{rwlang-cli} bináris, az alkalmazásforrás és a migration legyen ugyanahhoz a release-azonosítóhoz köthető. Ne egy élő \code{/srv/rwlang/current} könyvtár tartalmát módosítsd fájlonként; új immutable release könyvtárat készíts, majd atomikusan válts symlinket.

\section{systemd service}

A \code{rwlang-server} hosszú életű service process. Az \code{rwlang-cli} nem daemon: migration, auth-admin vagy más operátori művelet idejére indul és kilép.

\begin{lstlisting}
[Service]
User=rwlang
Group=rwlang
ExecStart=/usr/local/bin/rwlang-server \
  --config /usr/local/etc/rwlang/server.toml
Restart=on-failure
NoNewPrivileges=true
PrivateTmp=true
\end{lstlisting}

A service fusson unprivileged user alatt. A pontos systemd hardening és cgroup limitek igazodjanak a \code{server.toml} resource policy-jához; a két réteg célja ugyanaz: egy hibás vagy drága kérés ne vehesse át a gépet.

\section{Controlled restart és readiness}

Process-szintű konfigurációmódosítás után előbb fusson \code{--check-config}, majd controlled restart. Behind-proxy módban a SIGHUP a logok újranyitása mellett tranzakciósan újraépíti a domain/application hosting állapotot, de listener, DB/Redis/auth kapcsolat, cgroup policy, logging sink és más process-szintű beállítás továbbra is restartot igényel.

Az alkalmazás-forrás kiadásának van könnyebb útja. A közös source-reload supervisor minden betöltött alkalmazás entrypointját és teljes tranzitív modulgráfját \code{mtime + size} alapján figyeli. Konfigurálható debounce után candidate runtime-ot fordít és validál, majd csak siker esetén cseréli le atomikusan az adott domaint. Hibás feltöltésnél a korábbi generáció fut tovább, a reload-hiba pedig naplózódik. A logikai entrypoint útvonal is figyelt, ezért az immutable release-ek közötti atomikus \code{current} symlink-csere automatikusan aktiválja az új kódot. Sikeres kódreload a public-cache generationöket is előrelépteti, ezért a régi renderelt tartalom nem marad életben pusztán a korábbi TTL miatt.

A forgalomterelő csak ready instance-ra küldjön kérést. Rolling deploymentnél az új instance előbb induljon el, legyen ready, és csak utána kerüljön ki a régi. A health endpointok és a hozzájuk tartozó diagnosztika részletesen az observability fejezetben szerepel.

\section{Migration és rollback nem ugyanaz}

Az alkalmazás binárisának vagy forrásának visszaváltása egyszerű lehet, az adatbázis-séma visszaállítása viszont nem. Ezért production migrationnél preferáld az expand/contract mintát: az új release előbb backward-compatible sémát vezessen be, az inkompatibilis takarítás csak későbbi release-ben történjen.

Ne kapcsolj automatikus reverse SQL-t az alkalmazás symlink visszaállításához. Ha a migration adatot alakított át vagy törölt, a valódi recovery pont a bizonyított backup lehet.

\section{Recovery gate a deployment előtt}

Destruktív vagy schema-kompatibilitást érintő release előtt legyen ismert recovery point és friss restore-evidence. A backup tartalma, a DB/AppFs konzisztencia, az RPO/RTO, valamint a rollback és restore döntési fa a következő fejezetben szerepel. A deployment felelőssége itt az, hogy a release-váltás előtt ez a gate teljesüljön.

\section{Deploy utáni minimum ellenőrzés}

\begin{enumerate}
  \item az új process fut és ready;
  \item a publikus főoldal vagy smoke route válaszol;
  \item nincs új error spike a server logban;
  \item access logban a státuszkód-eloszlás normális;
  \item audit log írható;
  \item DB/Redis dependency metrikák nem romlottak;
  \item egy rollback döntéshez ismert a release hash és az adatbázis migration állapota.
\end{enumerate}
